% ============================================================
%  Team 13 — Medical Clinic Database Project
%  Compile with pdflatex (two passes for references/TOC)
% ============================================================
\documentclass[12pt, letterpaper]{article}

% ── Fonts & Encoding ────────────────────────────────────────
\usepackage[T1]{fontenc}
\usepackage[utf8]{inputenc}
\usepackage{lmodern}
\usepackage{microtype}

% ── Page Layout ─────────────────────────────────────────────
\usepackage[top=1in, bottom=1in, left=1.15in, right=1.15in]{geometry}
\usepackage{setspace}
\setstretch{1.15}

% ── Colors ──────────────────────────────────────────────────
\usepackage[dvipsnames, table]{xcolor}
\definecolor{clinicGreen}{HTML}{1f6d45}
\definecolor{clinicBlue}{HTML}{1f2a6d}
\definecolor{clinicTeal}{HTML}{2a7a5a}
\definecolor{lightGray}{HTML}{f5f5f5}
\definecolor{medGray}{HTML}{e0e3ed}
\definecolor{codeGray}{HTML}{f8f9fa}
\definecolor{codeBorder}{HTML}{d0d7de}
\definecolor{keywordBlue}{HTML}{0550ae}
\definecolor{commentGreen}{HTML}{1a7f37}
\definecolor{stringOrange}{HTML}{953800}
\definecolor{alertRed}{HTML}{cf222e}

% ── Graphics & Tables ───────────────────────────────────────
\usepackage{graphicx}
\usepackage{booktabs}
\usepackage{tabularx}
\usepackage{array}
\usepackage{multirow}
\usepackage{longtable}
\usepackage{colortbl}

% ── Lists & Enumeration ─────────────────────────────────────
\usepackage{enumitem}
\setlist[itemize]{leftmargin=1.4em, itemsep=2pt, topsep=4pt}
\setlist[enumerate]{leftmargin=1.4em, itemsep=2pt, topsep=4pt}

% ── Section Formatting ──────────────────────────────────────
\usepackage{titlesec}
\titleformat{\section}
  {\large\bfseries\color{clinicBlue}}
  {\thesection}{0.8em}{}
  [\color{medGray}\titlerule]
\titleformat{\subsection}
  {\normalsize\bfseries\color{clinicGreen}}
  {\thesubsection}{0.6em}{}
\titleformat{\subsubsection}
  {\normalsize\bfseries\color{clinicTeal}}
  {\thesubsubsection}{0.5em}{}
\titlespacing*{\section}{0pt}{18pt}{8pt}
\titlespacing*{\subsection}{0pt}{12pt}{4pt}
\titlespacing*{\subsubsection}{0pt}{8pt}{3pt}

% ── Code Listings ───────────────────────────────────────────
\usepackage{listings}
\lstset{
  basicstyle=\ttfamily\footnotesize,
  backgroundcolor=\color{codeGray},
  frame=single,
  framerule=0.4pt,
  rulecolor=\color{codeBorder},
  breaklines=true,
  breakatwhitespace=false,
  showstringspaces=false,
  tabsize=2,
  captionpos=b,
  aboveskip=8pt,
  belowskip=4pt,
  xleftmargin=8pt,
  xrightmargin=4pt,
  numbers=left,
  numberstyle=\tiny\color{gray},
  numbersep=6pt,
}

\lstdefinelanguage{MySQL}{
  keywords={SELECT, FROM, JOIN, LEFT, RIGHT, INNER, WHERE, AND, OR, NOT,
    ON, AS, ORDER, BY, GROUP, HAVING, LIMIT, INSERT, INTO, VALUES,
    UPDATE, SET, DELETE, CREATE, TABLE, TRIGGER, AFTER, BEFORE, FOR,
    EACH, ROW, BEGIN, END, IF, THEN, ELSE, CASE, WHEN, SIGNAL,
    SQLSTATE, MESSAGE, TEXT, DROP, EXISTS, DELIMITER, PROCEDURE,
    FUNCTION, RETURNS, DECLARE, CONCAT, ROUND, IFNULL, DATE, CURDATE,
    DATE\_ADD, DATE\_SUB, INTERVAL, NULLIF, SUM, COUNT, DISTINCT,
    NEW, OLD, NULL, NOT, IN, BETWEEN, LIKE, IS},
  keywordstyle=\color{keywordBlue}\bfseries,
  sensitive=false,
  comment=[l]{--},
  morecomment=[s]{/*}{*/},
  commentstyle=\color{commentGreen}\itshape,
  stringstyle=\color{stringOrange},
  morestring=[b]',
  morestring=[b]",
}

\lstdefinelanguage{JavaScript}{
  keywords={const, let, var, function, async, await, return, if, else,
    try, catch, throw, new, this, true, false, null, undefined,
    class, extends, export, import, from, of, for, while, do,
    switch, case, break, continue, typeof, instanceof},
  keywordstyle=\color{keywordBlue}\bfseries,
  sensitive=true,
  comment=[l]{//},
  morecomment=[s]{/*}{*/},
  commentstyle=\color{commentGreen}\itshape,
  stringstyle=\color{stringOrange},
  morestring=[b]',
  morestring=[b]",
  morestring=[b]`,
}

% ── Callout Box ─────────────────────────────────────────────
\usepackage{tcolorbox}
\tcbuselibrary{skins, breakable}

\newtcolorbox{infoblue}[1][]{
  enhanced, breakable,
  colback=blue!4!white, colframe=clinicBlue!60,
  fonttitle=\bfseries\small, title=#1,
  arc=4pt, left=8pt, right=8pt, top=6pt, bottom=6pt,
  before skip=8pt, after skip=6pt,
}

\newtcolorbox{greenbox}[1][]{
  enhanced, breakable,
  colback=green!4!white, colframe=clinicGreen!70,
  fonttitle=\bfseries\small, title=#1,
  arc=4pt, left=8pt, right=8pt, top=6pt, bottom=6pt,
  before skip=8pt, after skip=6pt,
}

\newtcolorbox{alertbox}[1][]{
  enhanced, breakable,
  colback=red!4!white, colframe=alertRed!60,
  fonttitle=\bfseries\small, title=#1,
  arc=4pt, left=8pt, right=8pt, top=6pt, bottom=6pt,
  before skip=8pt, after skip=6pt,
}

% ── Hyperlinks ──────────────────────────────────────────────
\usepackage{hyperref}
\hypersetup{
  colorlinks=true,
  linkcolor=clinicBlue,
  urlcolor=clinicTeal,
  pdftitle={Team 13 — Medical Clinic Database Project},
  pdfauthor={Team 13},
}

% ── Header / Footer ─────────────────────────────────────────
\usepackage{fancyhdr}
\pagestyle{fancy}
\fancyhf{}
\fancyhead[L]{\small\color{gray}Audit Trail Health — Team 13}
\fancyhead[R]{\small\color{gray}Medical Clinic Database}
\fancyfoot[C]{\small\thepage}
\renewcommand{\headrulewidth}{0.3pt}
\renewcommand{\footrulewidth}{0pt}

% ── Misc ────────────────────────────────────────────────────
\usepackage{parskip}
\setlength{\parskip}{6pt}
\usepackage{float}

% ============================================================
\begin{document}

% ── Title Page ──────────────────────────────────────────────
\begin{titlepage}
  \centering
  \vspace*{1.5cm}
  {\color{clinicBlue}\rule{\linewidth}{1.5pt}}\\[0.3cm]
  {\fontsize{26}{32}\selectfont\bfseries\color{clinicBlue}
    Audit Trail Health\\[0.15cm]}
  {\fontsize{14}{18}\selectfont\color{clinicGreen}
    Medical Clinic Database System}\\[0.3cm]
  {\color{clinicBlue}\rule{\linewidth}{1.5pt}}
  \vspace{1.2cm}

  {\large\bfseries Team 13}\\[0.4cm]
  {\normalsize
    Full-stack web application for a multi-clinic healthcare network.\\
    Patients, physicians, and staff each have role-specific portals.
  }

  \vspace{1.4cm}
  \begin{tabular}{ll}
    \textbf{Stack}    & Node.js + Express, MySQL (Railway), Vanilla JS \\
    \textbf{Branch}   & \texttt{TinaT2} \\
    \textbf{Database} & Railway MySQL (live demo) \\
    \textbf{Date}     & April 2026 \\
  \end{tabular}

  \vspace{1.8cm}
  \begin{greenbox}[Demo Login Credentials]
  \small
  \begin{tabular}{lll}
    \toprule
    \textbf{Role} & \textbf{Username} & \textbf{Password} \\
    \midrule
    Patient          & \texttt{alex.smith@email.com}    & \texttt{Patient@123} \\
    Patient          & \texttt{taylor.jones@email.com}  & \texttt{Patient@123} \\
    Physician (PCP)  & \texttt{dr.johnson}              & \texttt{Doctor@123} \\
    Physician (Spec) & \texttt{dr.white}                & \texttt{Doctor@123} \\
    Staff            & \texttt{staff.adams}             & \texttt{Staff@123} \\
    \bottomrule
  \end{tabular}
  \end{greenbox}

  \vfill
  {\small\color{gray}Compiled \today}
\end{titlepage}

\tableofcontents
\newpage

% ============================================================
\section{Project Overview}

Audit Trail Health is a full-stack web application simulating a real-world medical clinic network. It supports three distinct user roles — \textbf{patient}, \textbf{physician}, and \textbf{staff} — each with their own authenticated portal and specific data entry capabilities.

\subsection{Clinics \& Scale}
\begin{itemize}
  \item \textbf{3 clinic locations:} Dallas, Houston, and Austin, TX
  \item \textbf{65 physicians:} 38 primary care + 27 specialists across 8 specialties
  \item \textbf{Multiple departments} per clinic (primary care + specialist departments)
  \item \textbf{3 insurance plans:} BlueCross (80\%), Aetna (75\%), UnitedHealth (85\%)
\end{itemize}

\subsection{Technology Stack}
\begin{center}
\begin{tabular}{ll}
  \toprule
  \textbf{Layer}      & \textbf{Technology} \\
  \midrule
  Backend             & Node.js + Express.js \\
  Database            & MySQL via \texttt{mysql2} connection pool \\
  Hosting             & Railway (live demo database) \\
  Authentication      & \texttt{bcryptjs} password hashing (10 salt rounds) \\
  Frontend            & Vanilla HTML5 / CSS3 / JavaScript (no framework) \\
  Dev tooling         & Nodemon auto-restart, custom \texttt{npm run backup} script \\
  \bottomrule
\end{tabular}
\end{center}

% ============================================================
\section{Database Schema}

\subsection{Table Inventory}

\begin{longtable}{lll}
  \toprule
  \textbf{Table} & \textbf{Key Columns} & \textbf{Purpose} \\
  \midrule
  \endhead
  \texttt{clinic}             & clinic\_id, clinic\_name, city    & Top-level organization \\
  \texttt{department}         & dept\_id, clinic\_id, dept\_name  & Dept per clinic \\
  \texttt{office}             & office\_id, clinic\_id, city      & Physical location \\
  \texttt{physician}          & physician\_id, specialty, type    & All doctors \\
  \texttt{work\_schedule}     & physician\_id, day\_of\_week,     & Weekly recurring \\
                              & start/end\_time, office\_id       & availability \\
  \texttt{insurance}          & insurance\_id, coverage\_pct      & Insurance plans \\
  \texttt{staff}              & staff\_id, role, dept\_id         & Admin/nursing staff \\
  \texttt{users}              & user\_id, username, password\_hash, role & Login accounts \\
  \texttt{patient}            & patient\_id, user\_id, primary\_physician\_id & Patient profiles \\
  \texttt{appointment\_status}& status\_id, status\_name          & Lookup table \\
  \texttt{appointment}        & appt\_id, patient\_id, physician\_id, date, time & Visits \\
  \texttt{medical\_history}   & patient\_id, physician\_id, condition & Conditions/notes \\
  \texttt{diagnosis}          & appointment\_id, diagnosis\_code  & Per-visit diagnoses \\
  \texttt{treatment}          & diagnosis\_id, treatment\_plan, follow\_up\_date & Care plans \\
  \texttt{referral\_status}   & status\_id, status\_name          & Referral lookup \\
  \texttt{referral}           & patient\_id, primary\_physician\_id, specialist\_id & PCP→Specialist flow \\
  \texttt{billing}            & appointment\_id, total, ins\_paid, patient\_owed & Bills \\
  \bottomrule
\end{longtable}

\subsection{Key Relationships}
\begin{itemize}
  \item \texttt{users.physician\_id} $\rightarrow$ \texttt{physician.physician\_id} (physician login)
  \item \texttt{patient.primary\_physician\_id} $\rightarrow$ \texttt{physician.physician\_id} (nullable)
  \item \texttt{appointment} links patient, physician, office, and status in one row
  \item \texttt{billing} is auto-created by trigger when an appointment is marked Completed
  \item \texttt{referral}: \texttt{primary\_physician\_id} + \texttt{specialist\_id} both FK to \texttt{physician}
\end{itemize}

\subsection{Schema Enhancements Added}
The following columns were added beyond the base schema to support application logic:

\begin{itemize}
  \item \texttt{physician.physician\_type} \texttt{ENUM('primary','specialist')} — drives referral routing
  \item \texttt{medical\_history.physician\_id} — records which physician wrote each note
  \item \texttt{appointment.appointment\_type}, \texttt{appointment.duration\_minutes}
  \item \texttt{billing.insurance\_paid\_amount}, \texttt{billing.patient\_owed}, \texttt{billing.due\_date}
  \item \texttt{work\_schedule} — added \texttt{UNIQUE KEY(physician\_id, day\_of\_week)} constraint
\end{itemize}

% ============================================================
\section{Professor Requirements — Five Must-Haves}

% ── 1. Authentication ────────────────────────────────────────
\subsection{Requirement 1 — Authentication (All Three Roles)}

\begin{infoblue}[Status: Complete]
All three portals authenticate independently. Passwords are hashed with bcrypt (10 salt rounds). The database never stores plaintext.
\end{infoblue}

\subsubsection{How It Works}
\begin{enumerate}
  \item User submits email/username + password to \texttt{POST /api/auth/login}
  \item Server fetches the hashed password, calls \texttt{bcrypt.compare()}
  \item On success, returns \texttt{\{role, userId, patientId / physicianId / staffId\}}
  \item Frontend stores in \texttt{localStorage} and redirects based on role:
    \begin{itemize}
      \item \texttt{patient} $\rightarrow$ \texttt{/client/portals/patient\_dashboard.html}
      \item \texttt{physician} $\rightarrow$ \texttt{/client/portals/physician\_dashboard.html}
      \item \texttt{staff} $\rightarrow$ \texttt{/client/portals/staff\_dashboard.html}
    \end{itemize}
  \item Every subsequent API call passes \texttt{user\_id} and is validated by the \texttt{requireRole()} middleware
\end{enumerate}

\subsubsection{Security Features}
\begin{itemize}
  \item \textbf{Password strength enforcement:} 8+ chars, 1 uppercase, 1 number, 1 special character — validated both client-side and server-side
  \item \textbf{HIPAA idle auto-logout:} 15-minute inactivity timer on all three portals
  \item \textbf{Role isolation:} \texttt{requireRole()} middleware applied to every API route; a patient cannot call a physician endpoint and vice versa
  \item \textbf{Phone validation:} Registration requires a valid 10-digit US phone number (enforced frontend + backend)
  \item \textbf{Age validation:} Patients must be 18 or older; future birthdates rejected
\end{itemize}

% ── 2. Data Entry Forms ──────────────────────────────────────
\subsection{Requirement 2 — Data Entry Forms (Add, Modify, Delete)}

\begin{infoblue}[Status: Complete — all three portals]
Each role has forms covering add, modify, and delete operations.
\end{infoblue}

\subsubsection{Patient Portal}
\begin{center}
\begin{tabular}{lll}
  \toprule
  \textbf{Operation} & \textbf{Form} & \textbf{Route} \\
  \midrule
  ADD    & Book appointment (3-step modal)    & \texttt{POST /api/patient/appointments/book} \\
  MODIFY & Edit profile (name, phone, address) & \texttt{PUT /api/patient/profile} \\
  MODIFY & Assign/change care team            & \texttt{PUT /api/patient/care/assign} \\
  ADD    & Request specialist referral        & \texttt{POST /api/patient/referral/request} \\
  DELETE & Cancel appointment                 & \texttt{PUT /api/patient/appointments/:id/cancel} \\
  \bottomrule
\end{tabular}
\end{center}

\subsubsection{Physician Portal}
\begin{center}
\begin{tabular}{lll}
  \toprule
  \textbf{Operation} & \textbf{Form} & \textbf{Route} \\
  \midrule
  ADD    & Add clinical note to patient record   & \texttt{POST /api/staff/physician/note} \\
  MODIFY & Update appointment status             & \texttt{PUT /api/staff/appointment/:id/status} \\
  MODIFY & Review and action referrals (Issue/Reject) & \texttt{PUT /api/staff/referral/:id/status} \\
  MODIFY & Undo No-Show or Cancellation          & \texttt{PUT /api/staff/appointment/:id/undo-status} \\
  DELETE & Delete own medical history notes      & \texttt{DELETE /api/staff/medical-history/:id} \\
  \bottomrule
\end{tabular}
\end{center}

\subsubsection{Staff Portal}
\begin{center}
\begin{tabular}{lll}
  \toprule
  \textbf{Operation} & \textbf{Form} & \textbf{Route} \\
  \midrule
  ADD    & Book appointment for any patient & \texttt{POST /api/staff/appointments/book} \\
  MODIFY & Mark billing as paid            & \texttt{PUT /api/staff/billing/:id/pay} \\
  MODIFY & Update appointment status       & \texttt{PUT /api/staff/appointment/:id/status} \\
  \bottomrule
\end{tabular}
\end{center}

\subsubsection{Appointment Booking — 3-Step Form (Patient Portal)}
The booking modal is a guided 3-step data entry form:
\begin{enumerate}
  \item \textbf{Step 1 — Confirm Physician:} The patient's assigned primary physician is pre-loaded and locked. A hint explains that specialists require a referral.
  \item \textbf{Step 2 — Choose Date \& Time:} Displays the physician's available days of the week (e.g., ``Monday, Wednesday 9:00 AM–5:00 PM — Office: Dallas''). If the patient picks a non-working day, an inline error appears immediately without an API call. The date picker enforces a minimum of tomorrow and a maximum 2 years out.
  \item \textbf{Step 3 — Visit Details:} Appointment type and reason for visit. On submit, the double-booking trigger fires at the database level as a final safety check.
\end{enumerate}

% ── 3. Triggers ──────────────────────────────────────────────
\subsection{Requirement 3 — Database Triggers (3 implemented)}

\begin{greenbox}[Status: Complete — live on Railway]
All three triggers are deployed and active. They fire at the database layer regardless of application pathway.
\end{greenbox}

\subsubsection{Trigger 1 — Auto-Create Billing on Appointment Completion}

\textbf{Business rule:} When a physician marks an appointment Completed, the system automatically calculates and inserts a billing record using the patient's insurance coverage percentage. Zero manual data entry required.

\begin{lstlisting}[language=MySQL, caption={Trigger 1: after\_appointment\_completed}]
CREATE TRIGGER after_appointment_completed
AFTER UPDATE ON appointment
FOR EACH ROW
BEGIN
  IF NEW.status_id = 2 AND OLD.status_id != 2 THEN
    IF NOT EXISTS (
      SELECT 1 FROM billing WHERE appointment_id = NEW.appointment_id
    ) THEN

      SET @base_cost = CASE NEW.appointment_type
        WHEN 'Physical'   THEN 200.00
        WHEN 'Specialist' THEN 250.00
        WHEN 'Follow-Up'  THEN 125.00
        ELSE 150.00
      END;

      SET @coverage = IFNULL(
        (SELECT i.coverage_percentage FROM patient p
         JOIN insurance i ON p.insurance_id = i.insurance_id
         WHERE p.patient_id = NEW.patient_id), 0);

      SET @ins_id   = (SELECT insurance_id FROM patient
                       WHERE patient_id = NEW.patient_id);
      SET @ins_paid = ROUND(@base_cost * (@coverage / 100), 2);
      SET @owed     = ROUND(@base_cost - @ins_paid, 2);

      INSERT INTO billing (
        appointment_id, patient_id, insurance_id,
        total_amount, insurance_paid_amount, patient_owed,
        payment_status, due_date
      ) VALUES (
        NEW.appointment_id, NEW.patient_id, @ins_id,
        @base_cost, @ins_paid, @owed,
        'Unpaid', DATE_ADD(CURDATE(), INTERVAL 30 DAY)
      );
    END IF;
  END IF;
END
\end{lstlisting}

\textbf{Demo:} Physician marks an appointment Completed $\rightarrow$ switch to the billing tab $\rightarrow$ a fully calculated row appears with insurance math already applied.

\subsubsection{Trigger 2 — HIPAA No-Show Logging to Medical History}

\textbf{Business rule:} When an appointment is marked No-Show, an entry is automatically written to the patient's medical history with the date and visit context.

\begin{lstlisting}[language=MySQL, caption={Trigger 2: after\_appointment\_noshow}]
CREATE TRIGGER after_appointment_noshow
AFTER UPDATE ON appointment
FOR EACH ROW
BEGIN
  IF NEW.status_id = 4 AND OLD.status_id != 4 THEN
    IF NOT EXISTS (
      SELECT 1 FROM medical_history
      WHERE patient_id = NEW.patient_id
        AND `condition` = 'No-Show'
        AND diagnosis_date = NEW.appointment_date
    ) THEN
      INSERT INTO medical_history (
        patient_id, physician_id, `condition`,
        diagnosis_date, status, notes
      ) VALUES (
        NEW.patient_id,
        NEW.physician_id,
        'No-Show',
        NEW.appointment_date,
        'Active',
        CONCAT(
          'Patient did not attend scheduled appointment on ',
          DATE_FORMAT(NEW.appointment_date, '%M %d, %Y'),
          '. Reason for visit was: ',
          IFNULL(NEW.reason_for_visit, 'not specified'), '.'
        )
      );
    END IF;
  END IF;
END
\end{lstlisting}

\textbf{Demo:} Mark an appointment No-Show $\rightarrow$ medical history tab automatically shows a new entry with the appointment date and reason.

\subsubsection{Trigger 3 — Block Patient Double-Booking}

\textbf{Business rule:} A patient cannot have two appointments at the same date and time, even with different physicians. This fires before the insert and raises a database-level error forwarded to the UI.

\begin{lstlisting}[language=MySQL, caption={Trigger 3: before\_appointment\_double\_book}]
CREATE TRIGGER before_appointment_double_book
BEFORE INSERT ON appointment
FOR EACH ROW
BEGIN
  IF EXISTS (
    SELECT 1 FROM appointment
    WHERE patient_id       = NEW.patient_id
      AND appointment_date = NEW.appointment_date
      AND appointment_time = NEW.appointment_time
      AND status_id        != 3   -- ignore Cancelled slots
  ) THEN
    SIGNAL SQLSTATE '45000'
      SET MESSAGE_TEXT =
        'Patient already has an appointment at this date and time.';
  END IF;
END
\end{lstlisting}

\textbf{Demo:} Attempt to book the same patient at the same date and time $\rightarrow$ the booking modal displays the database error message directly.

% ── 4. Data Queries ──────────────────────────────────────────
\subsection{Requirement 4 — Data Queries (3 required + 2 bonus)}

\begin{greenbox}[Status: Complete — 5 queries implemented]
All queries use parameterized inputs (\texttt{?} placeholders) via \texttt{mysql2} to prevent SQL injection.
\end{greenbox}

\subsubsection{Query 1 — Patient Billing Statement}

Joins 5 tables to produce a complete per-appointment billing breakdown.

\begin{lstlisting}[language=MySQL, caption={Query 1: Patient Billing Statement}]
SELECT
  b.bill_id,
  a.appointment_date,
  a.appointment_type,
  CONCAT(ph.first_name, ' ', ph.last_name) AS physician_name,
  ph.specialty,
  o.city                                   AS office_city,
  b.total_amount,
  IFNULL(b.insurance_paid_amount, 0)       AS insurance_paid,
  IFNULL(b.patient_owed, 0)               AS patient_owed,
  b.payment_status,
  b.payment_method,
  b.payment_date,
  b.due_date,
  ins.provider_name,
  ins.coverage_percentage
FROM billing b
JOIN  appointment a   ON b.appointment_id = a.appointment_id
JOIN  physician ph    ON a.physician_id   = ph.physician_id
JOIN  office o        ON a.office_id      = o.office_id
LEFT JOIN insurance ins ON b.insurance_id = ins.insurance_id
WHERE b.patient_id = ?
ORDER BY a.appointment_date DESC;
\end{lstlisting}

\subsubsection{Query 2 — Daily Appointment Schedule}

Returns all appointments clinic-wide for a given date, ordered by location then time. Used by staff.

\begin{lstlisting}[language=MySQL, caption={Query 2: Daily Appointment Schedule}]
SELECT
  a.appointment_time,
  CONCAT(pt.first_name, ' ', pt.last_name) AS patient_name,
  pt.phone_number                          AS patient_phone,
  CONCAT(ph.first_name, ' ', ph.last_name) AS physician_name,
  ph.specialty,
  a.appointment_type,
  a.duration_minutes,
  a.reason_for_visit,
  s.status_name,
  o.city,
  o.street_address
FROM appointment a
JOIN patient pt              ON a.patient_id  = pt.patient_id
JOIN physician ph             ON a.physician_id = ph.physician_id
JOIN appointment_status s     ON a.status_id    = s.status_id
JOIN office o                 ON a.office_id    = o.office_id
WHERE a.appointment_date = ?
ORDER BY o.city, a.appointment_time;
\end{lstlisting}

\subsubsection{Query 3 — Physician Activity Report (Rolling 90-Day Window)}

Uses \texttt{CASE WHEN} inside \texttt{SUM} (conditional aggregation) to compute completion rate, revenue billed, no-show rate, and unique patients — all in a single query.

\begin{lstlisting}[language=MySQL, caption={Query 3: Physician Activity Report}]
SELECT
  ph.physician_id,
  CONCAT(ph.first_name, ' ', ph.last_name) AS physician_name,
  ph.specialty,
  COUNT(a.appointment_id)                  AS total_appointments,
  SUM(CASE WHEN s.status_name = 'Completed'
           THEN 1 ELSE 0 END)              AS completed,
  SUM(CASE WHEN s.status_name = 'No-Show'
           THEN 1 ELSE 0 END)              AS no_shows,
  SUM(CASE WHEN s.status_name = 'Cancelled'
           THEN 1 ELSE 0 END)              AS cancelled,
  ROUND(
    SUM(CASE WHEN s.status_name = 'Completed' THEN 1 ELSE 0 END)
    / NULLIF(COUNT(a.appointment_id), 0) * 100, 1
  )                                        AS completion_rate_pct,
  IFNULL(SUM(b.total_amount), 0)           AS total_revenue_billed,
  IFNULL(SUM(b.patient_owed), 0)           AS outstanding_balance,
  COUNT(DISTINCT a.patient_id)             AS unique_patients_seen
FROM physician ph
LEFT JOIN appointment a
  ON ph.physician_id = a.physician_id
  AND a.appointment_date
    BETWEEN DATE_SUB(CURDATE(), INTERVAL 90 DAY) AND CURDATE()
LEFT JOIN appointment_status s ON a.status_id    = s.status_id
LEFT JOIN billing b            ON a.appointment_id = b.appointment_id
WHERE ph.physician_id = ?
GROUP BY ph.physician_id, ph.first_name, ph.last_name, ph.specialty;
\end{lstlisting}

\subsubsection{Bonus Query 4 — Physician Utilization (Last 30 Days)}
Flags physicians with zero activity — useful for staff capacity planning.
\begin{itemize}
  \item Route: \texttt{GET /api/staff/physician-utilization}
  \item Uses \texttt{LEFT JOIN} to include physicians with no appointments
  \item Groups by physician, counts appointments, orders by activity ascending
\end{itemize}

\subsubsection{Bonus Query 5 — Billing Aging Buckets}
Groups unpaid bills into aging buckets: 0--30 / 31--60 / 61--90 / 90+ days overdue.
\begin{itemize}
  \item Route: \texttt{GET /api/staff/billing/aging}
  \item Uses \texttt{DATEDIFF(CURDATE(), due\_date)} with \texttt{CASE WHEN} to segment overdue bills
\end{itemize}

% ── 5. Data Reports ──────────────────────────────────────────
\subsection{Requirement 5 — Data Reports (3 implemented)}

\begin{greenbox}[Status: Complete — rendered in-app with print support]
Each report is generated on demand from a dedicated API route, rendered as a summary card + full data table, and includes a ``Print'' button (\texttt{window.print()}).
\end{greenbox}

\begin{center}
\begin{tabularx}{\linewidth}{llX}
  \toprule
  \textbf{Report} & \textbf{Portal} & \textbf{Shows} \\
  \midrule
  Billing Statement     & Patient $\rightarrow$ Billing tab    & Per-appointment: date, type, physician, total charged, insurance paid, patient owes, due date \\[4pt]
  Daily Schedule        & Staff $\rightarrow$ Appointments     & All clinic appointments for a selected date — time, patient, physician, type, reason, location, status \\[4pt]
  Physician Activity    & Physician $\rightarrow$ Reports tab  & 90-day rolling window: total appointments, completion rate \%, no-show rate, revenue billed, outstanding balance \\
  \bottomrule
\end{tabularx}
\end{center}

The reports differ from queries: queries return raw search results (interactive, with filters), while reports are pre-formatted outputs intended for review or printing.

% ============================================================
\section{Referral System}

The referral flow mirrors how specialist referrals work in the real US healthcare system:

\begin{center}
\begin{tabular}{cl}
  \toprule
  \textbf{Status} & \textbf{Meaning} \\
  \midrule
  Requested  & Patient requests a specialist via their PCP \\
  Issued     & PCP reviews and issues the referral to the specialist \\
  Accepted   & Specialist accepts the referral patient \\
  Rejected   & Specialist declines (patient must request again) \\
  Scheduled  & Patient books an appointment with the specialist \\
  Completed  & Specialist appointment is completed \\
  Expired    & Referral lapsed (90-day window) or patient changed PCP \\
  \bottomrule
\end{tabular}
\end{center}

\textbf{Business rules enforced:}
\begin{itemize}
  \item Patient must have at least one \textit{completed} appointment with their PCP before requesting a referral (eligibility gate)
  \item Only \texttt{physician\_type = 'primary'} physicians see the outgoing referral tab
  \item Only \texttt{physician\_type = 'specialist'} physicians see the incoming referral tab
  \item If a patient changes their PCP, all \textit{Requested} (not-yet-issued) referrals are automatically expired
  \item Referral cards on the patient portal show a ``Book Appointment with Dr. X'' button only when status is Accepted
\end{itemize}

% ============================================================
\section{Recent System Updates (Past 2 Days)}

The following features and fixes were implemented in the most recent development session.

\subsection{Booking Modal — Step Navigation \& Work Schedule Alignment}

\textbf{Bug fixed:} The modal was rendering all three steps simultaneously. Root cause: the \texttt{.hidden} CSS class was only defined for specific selectors (\texttt{.modal-overlay.hidden}, \texttt{.page-section.hidden}) but not generically. Added a global rule:

\begin{lstlisting}[language={},caption={CSS fix — global hidden utility}]
/* Global hidden utility — applies to any element with class="hidden" */
.hidden { display: none !important; }
\end{lstlisting}

\textbf{Work schedule alignment:} Before this fix, a patient could pick any calendar date in the booking modal and only discover a physician doesn't work that day after the API returned empty slots. Now:

\begin{enumerate}
  \item A new backend endpoint \texttt{GET /api/patient/appointments/physician-schedule} returns the physician's recurring weekly days and hours
  \item Step 2 of the booking modal loads and displays an availability card: \textit{``Physician availability: Monday 9:00 AM--5:00 PM, Wednesday 1:00 PM--5:00 PM — Office: Dallas''}
  \item When the patient selects a date, the weekday name appears below the input in real time
  \item If the selected day is not a working day, a red inline warning appears \textit{immediately} — no API call needed: \textit{``This physician doesn't work on Thursdays. Available days: Monday, Wednesday.''}
  \item Date input enforces \texttt{min} = tomorrow, \texttt{max} = 2 years out
\end{enumerate}

\subsection{Clinical Note Save Fix}

\textbf{Bug:} The ``Add Clinical Note'' form in the physician portal returned ``Failed to save note'' on every submission.

\textbf{Root cause:} The \texttt{POST /api/staff/physician/note} fetch call did not include \texttt{user\_id} in the request body, so the \texttt{requireRole()} middleware rejected every request with HTTP 401.

\textbf{Fix (two parts):}
\begin{itemize}
  \item Frontend: added \texttt{user\_id: user.id} to the fetch body in \texttt{submitNote()}
  \item Backend: \texttt{addPhysicianNote()} now looks up \texttt{physician\_id} from the \texttt{users} table and includes it in the \texttt{medical\_history} INSERT, so clinical notes are properly attributed to the physician who wrote them
\end{itemize}

\subsection{Appointment Date Display Fix}

\textbf{Bug:} Appointment dates rendered as \texttt{2026-11-11T00:00:00.000Z} in the physician portal Update button and other displays.

\textbf{Root cause:} JavaScript's \texttt{new Date("2026-11-11")} parses ISO date strings as UTC midnight, which then converts to the previous day in US local time zones (UTC-5/UTC-6).

\textbf{Fix:} Replaced the \texttt{fmt()} helper function throughout the physician portal:

\begin{lstlisting}[language=JavaScript, caption={Fixed date formatter — avoids UTC timezone shift}]
function fmt(d) {
    if (!d) return "—";
    const s = d.toString().split("T")[0];        // strip time component
    const [y, mo, dy] = s.split("-").map(Number);
    if (!y || !mo || !dy) return "—";
    return new Date(y, mo - 1, dy)               // parse as LOCAL date
        .toLocaleDateString("en-US", {
            month: "short", day: "numeric", year: "numeric"
        });
}
\end{lstlisting}

\subsection{Registration Phone Validation}

\textbf{Bug:} A single digit or empty phone number passed registration. The DOB validation block ran before the phone check, so an invalid DOB would stop the form before phone was ever evaluated.

\textbf{Fix:} Phone validation was moved \textit{before} DOB validation in \texttt{register.js}. A valid registration now requires:
\begin{itemize}
  \item Non-empty phone that normalizes to exactly 10 digits
  \item Auto-formatter displays as \texttt{(XXX) XXX-XXXX} while typing
  \item Backend (\texttt{authController.js}) applies the same rule server-side as a second layer
\end{itemize}

\subsection{Age-Aware Physician Recommendations}

In the patient care setup modal, the physician dropdown now shows contextual labels:
\begin{itemize}
  \item Patient is 65+, selecting a Geriatrics physician: \textit{``Dr. X — Geriatrics — Recommended for your age''}
  \item Patient under 65 sees Geriatrics: \textit{``--- Typically for patients 65+''}
  \item Family Medicine / Internal Medicine / General Practice: \textit{``--- All ages welcome''}
  \item Pediatrics: \textit{``--- For patients under 18''}
\end{itemize}

In the real world, Geriatrics is a valid primary care specialty — it is not a referral-only specialist. Geriatricians serve as primary care physicians for patients typically 65 and older.

\subsection{Physician Week-View Appointment Calendar}

A toggle between \textbf{List View} and \textbf{Calendar View} was added to the physician's Appointments section. The calendar:
\begin{itemize}
  \item Shows a 7-column week grid (Mon--Sun) with actual calendar dates in the column headers
  \item Today's column is highlighted in green
  \item Each appointment appears as a color-coded block (blue = Scheduled, green = Completed, amber = No-Show, gray = Cancelled) at the correct day
  \item Appointment blocks show: time, patient name, reason for visit
  \item Clicking any block opens the status update modal
  \item ``$\leftarrow$ Prev Week'' / ``Next Week $\rightarrow$'' navigation buttons shift the view
\end{itemize}

\subsection{Appointment Status Undo — With Justification Modal}

Physicians can now reverse a No-Show or Cancelled appointment back to Scheduled. This reflects real EMR behavior (Epic's ``Remove No-Show'' feature). Implementation details:

\begin{itemize}
  \item The appointments table shows an \textbf{Undo} button (amber) on No-Show and Cancelled rows; no Undo on Completed (billing already triggered)
  \item Clicking Undo opens a styled confirmation modal — not a browser \texttt{confirm()} dialog
  \item The modal shows patient name, appointment date, and current status
  \item A \textit{``Clinical Note''} panel explains the real-world context: reversing a No-Show requires justification in most EHR systems
  \item \textbf{Reason field is required} — the physician must type a justification before confirming
  \item On confirm: status reverts to Scheduled AND the reason is automatically logged as an \texttt{Appointment Status Correction} entry in the patient's \texttt{medical\_history} table — providing a full audit trail
\end{itemize}

\begin{lstlisting}[language=JavaScript, caption={Undo confirmation — reason required}]
async function confirmUndoStatus() {
    const reason = document.getElementById("undoReason").value.trim();
    if (!reason) {
        errEl.textContent =
          "Please enter a reason before restoring the appointment.";
        return;
    }
    const r = await fetch(
      `/api/staff/appointment/${_undoAppointmentId}/undo-status`,
      { method: "PUT",
        headers: { "Content-Type": "application/json" },
        body: JSON.stringify({ user_id: user.id, reason }) }
    );
    // On success: close modal, reload dashboard
}
\end{lstlisting}

\subsection{Care Team Change — Confirmation Box Alignment}

\textbf{Bug:} The red confirmation warning when changing physicians was rendered at full modal width while all other form fields had 28px horizontal padding, making it visually inconsistent.

\textbf{Fix:} Changed \texttt{margin-bottom:12px} to \texttt{margin:0 28px 12px} on the confirmation div so it aligns flush with the form fields above it.

\subsection{Staff Daily Schedule — Date Label}

The Daily Appointment Schedule modal header now shows the full date being viewed (e.g., \textit{``Daily Appointment Schedule — Today, Wednesday, April 15, 2026''}), making it clear which day's data is displayed when navigating to different dates.

% ============================================================
\section{Schema Gap Analysis}

The following known gaps are acknowledged. They do not affect the class demo but are worth noting for production use.

\begin{alertbox}[Known Schema Limitations]
\begin{enumerate}
  \item \textbf{work\_schedule has no date range} — \texttt{day\_of\_week} is a perpetual weekly pattern. There is no \texttt{effective\_from} / \texttt{effective\_to} and no \texttt{physician\_unavailable\_dates} table for holidays or vacation days. A physician scheduled for Mondays will theoretically show Monday slots every week indefinitely.
  \item \textbf{Insurance is plan-level, not patient-specific} — The \texttt{insurance} table stores plan-level coverage percentages. In a real system, each patient would have their own policy with a member ID, group number, and individual deductible.
  \item \textbf{No appointment capacity limits} — \texttt{work\_schedule} has no \texttt{max\_patients\_per\_day}. Slot generation is purely mathematical (30-minute intervals) and relies on the double-booking trigger for conflict prevention.
  \item \textbf{Referral expiry is passive} — The 90-day expiry is stored in \texttt{expiration\_date} but there is no scheduled job or trigger to automatically update status to ``Expired'' when the date passes.
  \item \textbf{No appointment history table} — Status changes are not tracked historically. The undo feature logs a note to \texttt{medical\_history} as a workaround, but a dedicated \texttt{appointment\_audit} table would be more appropriate.
\end{enumerate}
\end{alertbox}

% ============================================================
\section{API Route Reference}

\begin{longtable}{lll}
  \toprule
  \textbf{Method} & \textbf{Route} & \textbf{What it does} \\
  \midrule
  \endhead
  POST & \texttt{/api/auth/register}                        & Create user + patient row \\
  POST & \texttt{/api/auth/login}                           & Login, returns role \\
  GET  & \texttt{/api/patient/dashboard}                    & Patient dashboard data \\
  PUT  & \texttt{/api/patient/profile}                      & Save profile edits \\
  GET  & \texttt{/api/patient/appointments}                 & Patient's appointments \\
  GET  & \texttt{/api/patient/appointments/physician-schedule} & Physician working days \\
  GET  & \texttt{/api/patient/appointments/slots}           & Available time slots \\
  POST & \texttt{/api/patient/appointments/book}            & Book appointment \\
  PUT  & \texttt{/api/patient/appointments/:id/cancel}      & Cancel appointment \\
  PUT  & \texttt{/api/patient/care/assign}                  & Assign care team \\
  POST & \texttt{/api/patient/referral/request}             & Request referral \\
  GET  & \texttt{/api/staff/physician/dashboard}            & Physician dashboard \\
  POST & \texttt{/api/staff/physician/note}                 & Add clinical note \\
  PUT  & \texttt{/api/staff/appointment/:id/status}         & Update appt status \\
  PUT  & \texttt{/api/staff/appointment/:id/undo-status}    & Undo No-Show/Cancel \\
  PUT  & \texttt{/api/staff/referral/:id/status}            & Update referral status \\
  GET  & \texttt{/api/staff/dashboard}                      & Staff dashboard \\
  POST & \texttt{/api/staff/appointments/book}              & Staff books appointment \\
  PUT  & \texttt{/api/staff/billing/:id/pay}                & Mark billing paid \\
  GET  & \texttt{/api/reports/billing-statement}            & Report: billing \\
  GET  & \texttt{/api/reports/daily-schedule}               & Report: daily schedule \\
  GET  & \texttt{/api/reports/physician-activity}           & Report: activity \\
  \bottomrule
\end{longtable}

% ============================================================
\section{Presentation Notes}

\subsection{Talking Points for Triggers}
\begin{quote}
\itshape
``We implemented three AFTER/BEFORE triggers on the appointment table. Trigger 1 enforces a business rule: when a physician marks an appointment Completed, the system automatically calculates and inserts a billing record using the patient's insurance coverage percentage — zero manual data entry required. Trigger 2 implements HIPAA-aligned no-show tracking: when an appointment is marked No-Show, an entry is automatically written to the patient's medical history with the date and visit context. Trigger 3 is a BEFORE INSERT guard that prevents double-booking at the database level, raising a SQL error that surfaces directly to the booking UI. All three triggers fire regardless of which application pathway updates the record.''
\end{quote}

\subsection{Talking Points for Queries}
\begin{quote}
\itshape
``We wrote five parameterized SQL queries spanning multiple joined tables. Query 1 joins five tables to produce a patient billing statement covering every appointment, its physician, office, insurance coverage, and outstanding balance. Query 2 generates a daily schedule across all locations, useful for clinic operations. Query 3 is an analytical query using conditional aggregation — CASE WHEN inside SUM — to compute a physician's completion rate, revenue, and no-show rate over a rolling 90-day window. This kind of query is standard in healthcare analytics.''
\end{quote}

\subsection{Talking Points for Reports vs.\ Queries}
\begin{quote}
\itshape
``Reports are pre-formatted, printable outputs generated on demand. Queries are interactive — a staff member types a search and gets live filtered results. Both are implemented differently and both count toward the professor's requirements.''
\end{quote}

\end{document}
